\documentclass[12pt]{article}

\usepackage{amsmath}
\usepackage{german}
\usepackage{graphicx}
\DeclareMathOperator\sgn{sgn}

%\graphicspath{{/home/ritz/Teaching/fcn/fig/}}

%----------------- layout ----------------------------------------

\oddsidemargin=0in
\textwidth=6in
\topmargin=-0.7in
\textheight=9.5in
\baselineskip=16pt
\font\mittel=cmr17
\font\roman=cmr12
\font\klein=cmr9
\font\mgross=cmr10
\newfont{\CAP}{cmcsc10 scaled 900}
\pagestyle{empty}
\parindent=0mm


%----------------- macros ----------------------------------------
\newcommand{\conc}[1]{[{\rm C}]_{\text{#1}}}
\newcommand{\units}[1]{\mathrm{#1}}
\newcommand{\K}{\text{K}^+}
\newcommand{\Ca}{\text{Ca}^{2+}}
\newcommand{\Na}{$\text{Na}^+$}
\newcommand{\Cl}{$\text{Cl}^-$}
\newcommand{\Leak}{\text{leak}}
\newcommand{\water}{$\text{H}_2\text{O}$}
%-----------------------------------------------------------------


\begin{document}

\parbox{2cm}{
\includegraphics[width=1.8cm]{bccnlogo.pdf}
}
\parbox{11cm}{
\begin{center}
\large HUMBOLDT-UNIVERSIT"AT \hskip 0.1 cm ZU \hskip 0.1 cm BERLIN
\vskip 0.1 true cm
\mgross BERNSTEIN CENTER FOR\\ COMPUTATIONAL NEUROSCIENCE
\end{center}
}
\parbox{2cm}
{
\hfill
\includegraphics[width=1.8cm]{hublogo.pdf}
}

\vskip 0.6 true cm

\leftline{\CAP Humboldt-Universit"at zu Berlin 
\hfill Phone: 030/2093-9110}
\leftline{\CAP Philippstr. 13 House 6
              \hfill Fax: 030/2093-6771}
\leftline{\CAP \hfill webpage: http://www.bccn-berlin.de/}


\vskip 0.8 true cm
\centerline{\bf Models of Neural Systems, WS 2010/11}
\centerline{\bf Project 8: Bursting}
\centerline{Project presentation and report submission: February, 15th, 2011}

\vskip 0.8 true cm
\noindent {\bf Background}

Many neurons can produce brief trains of spikes interspersed with
long quiescent periods. This phenomenon called bursting was found in many
brain areas such as thalamus and cortex and may play an important role
in information transmission and encoding.

The aim of the project is to study burst generation in
pancreatic $\beta$-cells. The membrane dynamics of the cells can be
described by a simplified system of 3 differential equations: 
\begin{eqnarray}
    C_m\frac{dV}{dt}&=&-I_\Ca(V)-\left(\bar{g}_\K
    n^4+\frac{\bar{g}_{\K,\Ca}c}{K_d+c}\right)(V-V_\K)-\bar{g}_\Leak(V-V_\Leak),
    \label{eq:pancreatic1}\\
\tau_n(V)\frac{dn}{dt}&=&n_\infty(V)-n,\label{eq:pancreatic2} \\
\frac{dc}{dt}&=&f(-k_1 I_\Ca(V)-k_cc), \label{eq:pancreatic3}
\end{eqnarray}

where $V$ is a membrane potential, $n$ is a Hodgkin-Huxley type
activation variable  and $c$ is a slow variable
describing the intracellular concentration of calcium.

  The $\Ca$ current flows through a voltage-gated channel and thus:
\begin{equation}
I_\Ca(V)=\bar{g}_\Ca m_\infty^3(V)h_\infty(V)(V-V_\Ca)
\end{equation}

The steady state activation and time constant of $m$ and $h$ variables
are the same as in the Hodgkin-Huxley model, except that the voltage is
shifted by $V^\star=-15\,\units{mV}$ so that $V$ is replaced by
$V+V^\star$. Similarly, $n$ satisfies Hodgkin-Huxley type differential equation, shifted along the voltage axis by an amount $V^+=-35\,\units{mV}$.

The remaining parameters of the model are: $C_m=1\,\units{\mu
F/cm^2}$, $\bar{g}_\K=5.5\,\units{mS/cm^2}$,
$\bar{g}_\Leak=0.03\,\mathrm{mS/cm^2}$, $V_\Ca=100\,\units{mV}$,
$K_d=1\,\units{\mu M}$, $k_1=0.0275\,\units{\mu M\:cm^2/nC}$,
$\bar{g}_{\K,\Ca}=0.05\,\units{mS/cm^2}$,
$\bar{g}_\Ca=8\,\units{mS/cm^2}$, $V_\K=-75\,\units{mV}$,
$V_\Leak=-40\,\units{mV}$, $f=0.007$, $k_c=0.02\,\units{ms^{-1}}$.

Use the initial conditions $c(t=0)=0.45$, $V(t=0)=-50\,\units{mV}$, $n(t=0)=0$.

\vskip 1.7 true cm

{\bf Problems}


\begin{enumerate}
    \item Simulate the model of a pancreatic $\beta$-cell given
        by equations (\ref{eq:pancreatic1})-(\ref{eq:pancreatic3}).
        Show the bursting phenomenon.
    \item Plot the histogram of interspike intervals (ISI). Find the modes
        of the histogram relating to within burst intervals (WBI) and
        interburst intervals (IBI). Plot the WBI lengths as the
        function of the spike position in the burst.
    \item Treating $c$ as a constant parameter, consider the
        phase-plane structure of fast
        subsystem consisting of equations
        (\ref{eq:pancreatic1})-(\ref{eq:pancreatic2}). Try to find the
        values of parameter $c$ for which a stable point and a limit
        cycle exists.
    \item Plot the bifurcation diagram of the fast subsystem
        varying parameter $c$. Plot the stable/unstable fixed points
        and stable limit cycles. Identify the bifurcation points. 
    \item Plot the nullcline of the slow variable on the bifurcation
        diagram. Can you explain the mechanism of the bursting using
        the diagram? Superimpose a $(V,c)$-projection of a burst cycle on the plot.
    \item Using the bifurcation diagram predict which values of the $k_c$
        parameter will lead to bursting. 
\end{enumerate}

\begin{thebibliography}{7}
    \bibitem{keener2008} J Keener, J. Sneyd,``Mathematical physiology'', Chapter 6: Bursting Electrical Activity
    \bibitem{chay1983} TR Chay, J Keizer, ``Minimal model for membrane oscillations in the pancreatic beta-cell``, Biophysical Journal 42 (2): 181-189 (1983)
\end{thebibliography}

\vfill
\centerline{\CAP Contact}
\CAP

\begin{tabular}{lll}
\underline{Matthias Guggenmos} & Phone: 450-517131 & Email: matthias.guggenmos(at)bccn-berlin.de \\
Urs Bergmann & Phone: 2093-8924 & Email: urs.bergmann(at)biologie.hu-berlin.de \\
Bartosz Telenczuk & Phone: 2093-8838 & Email: b.telenczuk(at)biologie.hu-berlin.de \\
Richard Kempter & Phone: 2093-8925 & Email: r.kempter(at)biologie.hu-berlin.de \\
\end{tabular}
\end{document}

 
